\everymath{\displaystyle}

%%%%%%%%%%%%%%%%%%%%%%%%%%%%%%%%%%%%%%%%
%           main body
%%%%%%%%%%%%%%%%%%%%%%%%%%%%%%%%%%%%%%%%

%-----------------------------------
%	TITLE PAGE
%-----------------------------------

\title[第七章\\
定积分]{\texorpdfstring{\LARGE \bfseries 第七章\quad 定积分}{第七章 定积分}} % The short title appears at the bottom of every slide, the full title is only on the title page
\author[\smaller \S 7.3\\
微积分基本定理] % Your institution as it will appear on the bottom of every slide, maybe shorthand to save space
{\Large \bfseries \S 7.3 微积分基本定理}
% \author{Singular Value Decomposition} % Your name
% \date{\today} % Date, can be changed to a custom date
\date{}

%------------------------------------------------

\begin{document}

%------------------------------------------------

\begin{frame}

\titlepage % Print the title page as the first slide

\end{frame}

%------------------------------------------------

% \begin{frame}
% \frametitle{内容提要} % Table of contents slide, comment this block out to remove it
% \tableofcontents % Throughout your presentation, if you choose to use \section{} and \subsection{} commands, these will automatically be printed on this slide as an overview of your presentation
% \end{frame}

%-----------------------------------
%	PRESENTATION SLIDES
%-----------------------------------

%------------------------------------------------

\begin{frame}
\frametitle{绪言}

{
\larger[2]

\begin{center}
这节要将\underline{积分学}和\underline{微分学}联系起来.
\end{center}

\pause
\vspace{1em}

之前已经接触过 $F(x) = \int_a^x f(t) ~ \mathrm{d} t$ 这个函数了, 这里 $f \in \mathbb{R}[a, b]$ 是一个黎曼可积函数. $F(x)$ 是所谓的\alert{变上限积分},

}

\end{frame}

%------------------------------------------------

\begin{frame}
\frametitle{待写}

\end{frame}

%------------------------------------------------

\end{document}
